\worksheettitle
  {Домашняя работа: составляем числа}
  {\modulemark{M05-H} Самостоятельное закрепление}

\studentnameblock

\begin{notebox}[Перед началом]
  Для каждого задания отметьте: разрешённые цифры, длину числа, правило
  повторений и обязательность цифр. Списки располагайте по возрастанию первой
  цифры.
\end{notebox}

\begin{practicebox}[1. Составьте списки]
  \begin{enumerate}[label=\textbf{\arabic*.}]
    \item Из цифр \(2\) и \(6\) составьте все двузначные числа. Повторения
      разрешены, использовать обе цифры необязательно:
      \answerblank[72mm].
    \item Из цифр \(2,\ 6,\ 8\) составьте все двузначные числа без повторений:
      \answerblank[84mm].
    \item Из цифр \(0,\ 3,\ 9\) составьте все трёхзначные числа, используя
      каждую цифру ровно один раз:
      \answerblank[84mm].
  \end{enumerate}
\end{practicebox}

\begin{practicebox}[2. Обязательные цифры]
  Из цифр \(4\) и \(7\) составьте все трёхзначные числа. Повторения разрешены,
  обе цифры обязательны:
  \answerblank[98mm].

  Какие два числа не вошли и почему?
  \writinglines{1}
\end{practicebox}

\begin{practicebox}[3. Докажите полноту]
  Выберите один список из задания 1. Объясните, почему в нём нет пропусков и
  повторов.
  \writinglines{3}
\end{practicebox}

\newpage
\begin{warningbox}[4. Исправьте чужой список]
  Из цифр \(1,\ 4,\ 8\) составляли двузначные числа без повторений:
  \(14,\ 18,\ 41,\ 48,\ 81,\ 81\).

  Повтор: \answerblank[15mm]. Пропущено: \answerblank[15mm].
  Исправленный список: \answerblank[82mm].
\end{warningbox}

\begin{warningbox}[5. Ноль в начале]
  Для цифр \(0,\ 5,\ 8\) без повторений записали:
  \(05,\ 08,\ 50,\ 58,\ 80,\ 85\).

  Какие записи не являются двузначными числами? \answerblank[30mm].
  Полный список двузначных чисел: \answerblank[70mm].
\end{warningbox}

\begin{practicebox}[6. Составьте условие]
  Для списка \(22,\ 25,\ 52,\ 55\) сформулируйте точное условие задачи.
  \writinglines{2}

  Как изменится список, если обе цифры обязательны?
  \answerblank[50mm].
\end{practicebox}

\begin{checkpointbox}[7. Итог]
  \begin{enumerate}[label=\textbf{\arabic*.}]
    \item Из цифр \(0,\ 2,\ 7\) составьте все трёхзначные числа, используя
      каждую цифру ровно один раз:
      \answerblank[82mm].
    \item Объясните порядок списка и отсутствие вариантов с первой цифрой 0.
      \writinglines{2}
  \end{enumerate}
  \textbf{Перед сдачей:} проверьте количество вариантов и отсутствие повторов.
\end{checkpointbox}
